% ======================================================================
%  CAPÍTULO 18 — SOLUCIONES (parte 3: capítulos 9-16)
% ======================================================================

\subsection*{Capítulo 9 --- La generación de claves
(capítulo~\ref{ch:keygen})}

\begin{enumerate}[leftmargin=1.9em,itemsep=4pt]
  \item Todo el proceso, incluidos los reintentos, es una función
        determinista de la semilla $\delta$: al fallar
        $\textsf{FieldOrdering}$, se toma $\delta'\leftarrow
        \textnormal{(últimos $256$ bits de }E)$ y se repite sin pedir
        azar nuevo. Las pruebas de vectores fijan $\delta$ y el
        resultado esperado: cada implementación consume el mismo flujo
        $E$ y reproduce los mismos reintentos, comparables byte a
        byte.
  \item Con $(7,1,9,3)$: al ordenar los pares $(a_i,i)$ queda
        $1<3<7<9$, es decir $\pi=(1,3,0,2)$; leyendo los bits de
        $\pi(i)$ como $\alpha_i=\sum_{j=0}^{m-1}\pi(i)_jz^{m-1-j}$
        resulta $(\alpha_0,\dots,\alpha_3)=(z,\,z+1,\,0,\,1)$. Con
        $(0,1,2,3)$: $\pi$ es la identidad
        y $(\alpha_0,\dots,\alpha_3)=(0,\,z,\,1,\,z+1)$. Con
        $(0,0,0,0)$ los enteros no son distintos:
        $\textsf{FieldOrdering}$ devuelve $\bot$ y
        $\textsf{MatGen}$ reintenta con $\delta\leftarrow\delta'$
        (determinista).
  \item $\beta$ recorre uniformemente $\GF{q^t}$ (los bloques de
        $\sigma_1=16$ bits por coeficiente; solo se usan los primeros
        $m$ de cada bloque). El polinomio mínimo de $\beta$ sobre
        $\GF q$ es mónico, irreducible, y de grado $t$ exactamente
        cuando $\beta$ no cae en un subcuerpo propio
        $\GF{q^s}$, $s\mid t$, $s<t$ (si cae, el grado es $s$ y
        $\textsf{Irreducible}$ devuelve $\bot$). Cada mónico
        irreducible de grado $t$ divide a $x^{q^t}-x$ y tiene
        exactamente $t$ raíces, todas simples, en $\GF{q^t}$; sus
        $t$ raíces son exactamente las $\beta$ cuyo mínimo es él. Así
        cada $g$ admite $t$ de las $q^t$ elecciones: $\beta$
        uniforme $\Rightarrow$ $g$ uniforme (la masa rechazada es
        $\approx q^{t-1}/q^t=1/q$).
  \item La probabilidad de que las primeras $mt$ columnas de una
        matriz binaria aleatoria $mt\times n$ sean independientes es
        $\prod_{i=0}^{mt-1}(1-2^{i-mt})
        =\prod_{j=1}^{mt}(1-2^{-j})\approx0{,}289$:
        $3{,}46$ intentos de media, y $P(\ge3)\approx(1-0{,}289)^2
        \approx0{,}5055$. Con $(\mu,\nu)=(32,64)$ el fallo por
        intento cae bajo $2^{-30}$: un solo intento en la práctica,
        a cambio de $\le2$ bits de tensión en la reducción y una
        ventana de intercambios de tamaño fijo (tiempo constante).
  \item $pk=mt\lceil k/8\rceil=1664\cdot816=\num{1357824}$~B;
        $sk=32+8+t\lceil m/8\rceil+\lceil(2m-1)2^{m-4}\rceil
        +\lceil n/8\rceil=40+256+12\,800+1024=\num{14120}$~B.
        En la clave secreta domina $\alpha$ (los $12\,800$ bytes de
        la red de Benes, el $91\%$); en la pública, la matriz $T$
        entera.
  \item El código $\ker(\mathrm{I}_{mt}|T)$ es público: dimensión
        $k=2720$ y $2^{2720}$ codewords en \texttt{348864}
        (cualquiera con la clave pública comprueba membresías). Quien
        descubra $g$ y el orden $\alpha$ obtiene la descripción de
        Goppa del código público: un decodificador eficiente, es
        decir, una clave equivalente: descifrar todo
        (observación~\ref{obs:code-public} y
        proposición~\ref{prop:nied-reduccion}: atacar el esquema es
        decodificar el código público; el secreto es su
        descripción).
\end{enumerate}

\subsection*{Capítulo 10 --- Cifrado y descifrado
(capítulo~\ref{ch:encapdecap})}

\begin{enumerate}[leftmargin=1.9em,itemsep=4pt]
  \item $v=(C,\mathbf{0})$; como $C=\hat{\Hmat}e$, para $c=v+e$
        vale $\hat{\Hmat}c=\hat{\Hmat}v+\hat{\Hmat}e=C+C=0$: es un
        codeword, y $d(c,v)=\wt(e)=2=t$. Con $C'=\hat{\Hmat}e'$ y
        $\wt(e')=1$: el paso~2 de $\textsf{Decode}$ encuentra el
        único codeword $c'$ a distancia $\le2$ (que es $v' + e'$);
        el paso~3 devuelve $e''=e'$ (peso $1$); el paso~4
        \emph{rechaza} ($\wt(e')\ne t$) y $\textsf{Decap}$ toma la
        rama de rechazo.
  \item (a)~Los trozos de $16$ bits bastan para nombrar cualquier
        posición porque todo índice es $<n\le q=2^m\le2^{16}$ (en
        los conjuntos reales $m\le13$; el juguete usa $m=4$);
        (b)~cada
        elección usa un trozo fresco: las elecciones son
        independientes; (c)~descartar los trozos fuera de rango y
        repetir no sesga: condicionado a ser aceptado, cada residuo
        es equiprobable. Productos: juguete:
        $1-1/16=0{,}9375$; \texttt{348864}:
        $e^{-t(t-1)/(2n)}=e^{-64\cdot63/6976}\approx0{,}561$.
  \item Si $c_1\ne c_2$ estuvieran ambos a distancia $\le t$ de
        $v$, entonces $\wt(c_1+c_2)\le2t<2t+1\le d$: la diferencia
        sería un codeword no nulo de peso bajo la distancia mínima.
        Imposible: el candidato es único.
  \item Con $C'=\hat{\Hmat}e'$, $\wt(e')=1$: el paso~2 devuelve
        $e'$ (la corrección funciona: $1\le t$); el paso~4 lo
        rechaza por el chequeo de peso. Sin ese chequeo,
        $\textsf{Decap}$ aceptaría todo síndrome con preimagen de
        peso $\le t$, y los «mensajes» aceptados serían uniformes en
        $\{$peso $\le t\}$ en vez de en $\{$peso $=t\}$: se
        romperían la distribución y la rigidez que las pruebas del
        capítulo~\ref{ch:kem} usan (la aceptación debe caracterizar
        exactamente los criptogramas de $\textsf{FixedWeight}$).
  \item \textsf{Encode}: $mt\cdot k=1664\cdot6528=\num{10862592}$
        bits $\approx2^{23{,}4}$; síndrome de \textsf{Decode}:
        $tn=128\cdot8192=\num{1048576}$ $\approx2^{20}$ operaciones
        de cuerpo. \textsf{Encode} toca unos $10$ veces más bits,
        pero son XORs sobre datos públicos (vectorizables, sin
        secretos); las de \textsf{Decode} son multiplicaciones en
        $\GF{2^{13}}$, más Patterson ($t^2$ operaciones) y la
        verificación: en coste efectivo domina \textsf{Decode}. Y
        esa asimetría es justo la que quiere un KEM: el lado sin
        clave (cifrar) por el camino barato y paralelo, y el lado
        con la clave (descifrar) pagando el álgebra fuerte.
  \item Para \texttt{6960119}: $H(b,v,C)$ recibe $1+\lceil
        n/8\rceil+\lceil mt/8\rceil=1+870+194=1065$ bytes. Piezas:
        la etiqueta ($0/1/2$) separa dominios (rechazo, éxito,
        confirmación \texttt{pc}); $v=e$ fija el error; $C$ ata la
        clave a la codificación exacta del criptograma (el escenario
        del capítulo~\ref{ch:kem}). No se usa la lista de posiciones
        del soporte porque el vector ya es una codificación de
        longitud fija y canónica ($\lceil n/8\rceil$ bytes): una
        lista exigiría convenios de orden, tendría longitud variable
        y no aportaría nada.
\end{enumerate}

\subsection*{Capítulo 11 --- Laboratorio 2
(capítulo~\ref{ch:labkem})}

\begin{enumerate}[leftmargin=1.9em,itemsep=4pt]
  \item Ejecutando el programa del apéndice~\ref{app:codigo}:
        $\textsf{SHAKE256}(\hx{01}\,\|\,\hx{00\,24}\,\|\,\hx{98})$
        da exactamente la $K$ de la
        figura~\ref{fig:labkem-pipeline}
        (\texttt{597f\ldots6278}), y $s=\hx{CBAE}$ se serializa
        como los bytes $\hx{AE}\,\hx{CB}$ (little-endian por grupos
        de $8$).
  \item $e=9216=2^{10}+2^{13}$: las posiciones $10$ y $13$; cada
        byte empaqueta $8$ posiciones con la posición $j$ en el bit
        $j\bmod8$ del byte $j/8$: el primer byte
        (posiciones $0$--$7$) es $\hx{00}$ y el segundo
        (posiciones $8$--$15$) es $2^2+2^5=\hx{24}$: los bytes
        $\hx{00}\,\hx{24}$. La convención «posición $0$ a la
        izquierda» describe cómo se \emph{imprimen} los vectores;
        la serialización empaqueta en little-endian. Igual para
        $s$: $\hx{CBAE}$ escribe $\hx{AE}\,\hx{CB}$.
  \item Si $c=v+e$ es un codeword y $S_c=\sum_jc_j\sigma_j\equiv0
        \pmod g$, entonces
        $S_v=\sum_jv_j\sigma_j\equiv S_c+S_e\equiv S_e\pmod g$
        (en $\Ftwo$, $v_j+e_j=c_j$). Para $\tilde C$ el vector
        $\tilde v+e$ ya \emph{no} es un codeword: su síndrome es la
        columna del bit volteado ($\ne0$), la cancelación falla y el
        síndrome cambia (en el juguete, $S=\hx4+\hx9x$ pasa a
        $\hx b+\hx6x$).
  \item Si $\textsf{Decap}$ dijera «error», el atacante voltearía
        bits en $C$ y aprendería de la respuesta si $\tilde C$
        decodifica: un oráculo de validez sobre los $256$ bytes
        (los $119$ inalcanzables marcan los que siempre fallan; cada
        consulta revelaría clases). En el KEM real no hay señal: la
        rama de rechazo devuelve $H(0,s,\tilde C)$, indistinguible
        de una $K$ válida; los $119$ síndromes inalcanzables se
        tratan como cualquier otro fracaso.
  \item $1+16+120=137$ síndromes alcanzables con peso $\le2$; los de
        peso $3$ ya no están garantizados como únicos porque la
        diferencia de dos palabras de peso $3$ puede pesar hasta
        $6>d=5$ (solo los pesos con sumas $<d$ son únicos).
        $137/256\approx53{,}5\%$ es la fracción de bytes alcanzables
        por algún error de peso $\le2$; un byte manipulado al azar
        falla tras el chequeo de peso con probabilidad
        $(256-120)/256=136/256\approx53{,}1\%$: casi la misma cifra
        (los $17$ síndromes de peso $0/1$, rechazados por el
        chequeo, reparten la diferencia).
  \item Clave pública: $\num{261120}/8=\num{32640}$ veces la del
        juguete; criptograma: $96/1=96$ veces. Mecanismos omitidos:
        $\sigma_1=16$ (bits sobrantes por coeficiente en
        $\textsf{Irreducible}$), $\sigma_2=32$ (bits por elemento en
        $\textsf{FieldOrdering}$), la forma semi-sistemática
        $(\mu,\nu)$ y la red de Benes (compresión del orden
        $\alpha$): los cuatro, en el capítulo~\ref{ch:keygen} (y el
        tercero también en el~\ref{ch:parametros}).
\end{enumerate}

\subsection*{Capítulo 12 --- Los conjuntos y las variantes
(capítulo~\ref{ch:parametros})}

\begin{enumerate}[leftmargin=1.9em,itemsep=4pt]
  \item Con $ct=\lceil mt/8\rceil$ y $ek=mt\lceil k/8\rceil$:
        \texttt{348864}: $96$ y $768\cdot340=\num{261120}$;
        \texttt{460896}: $156$ y $1248\cdot420=\num{524160}$;
        \texttt{6688128}: $208$ y $1664\cdot628=\num{1044992}$;
        \texttt{6960119}: $194$ y $1547\cdot677=\num{1047319}$;
        \texttt{8192128}: $208$ y $1664\cdot816=\num{1357824}$.
        En todos menos \texttt{6960119}, $k$ es múltiplo de $8$ y
        $mt\lceil k/8\rceil=\lceil kmt/8\rceil$. En
        \texttt{6960119}, $k=5413$: $\lceil kmt/8\rceil
        =\lceil\num{8373911}/8\rceil=\num{1046739}$, frente a
        $\num{1047319}$: la fila se redondea a byte y se pierden
        $580$~B (el relleno de la última fila, $3$ bits por fila).
  \item Con $H(2,e)$: $1+\lceil n/8\rceil$ bytes;
        $H(1,e,C)$ y $H(0,s,C)$: $1+\lceil n/8\rceil+\lceil mt/8\rceil
        +32$ en las variantes \texttt{pc} (el $32$ es la
        confirmación). Para \texttt{460896pc}: $577$ y $765$~B; la
        serie: $437/577/837/871/1025$ (para $H(2,e)$) y, en
        \texttt{pc}, $565/765/997/1097/1265$ (para los otros dos).
  \item Mitades de Stern: $75{,}7$, $96{,}7$, $134{,}0$,
        $134{,}2$, $151{,}6$. Superan $2^{128}$ los tres conjuntos
        $\texttt{6*}$: exactamente el borrador ISO de 2023. ISO
        estandarizó \emph{además} \texttt{460896}, sabiendo que no
        alcanza el listón ingenuo, porque la raíz cuadrada es
        optimista para el atacante (Grover limitada por latencia):
        el borrador lo dejó fuera y la versión final lo incluyó con
        esa salvedad (la historia del capítulo~\ref{ch:parametros}).
  \item $p\approx0{,}289$: $3{,}46$ intentos de media y
        $P(>2)=(1-p)^2\approx0{,}5055$. Con $(\mu,\nu)=(32,64)$ el
        fallo es $<2^{-30}$: $\approx1$ intento. La variante
        \texttt{f} se ahorra $\approx3{,}5$ generaciones por clave
        (y elimina la cola y la varianza).
  \item Dos cadenas de $194$ bytes que difieran solo en los $5$
        bits de relleno representan el mismo elemento de
        $\Ftwo^{1547}$. En modo Simply el relleno se ignora: ambas
        producen la \emph{misma} $K$ (un solo valor distinto). En
        modo Narrowly, la cadena con relleno no nulo se rechaza:
        una produce la $K$ legítima y la otra la clave de la rama de
        rechazo. Debe fijarlo el protocolo porque una implementación
        Simply y una Narrowly, mezcladas, discreparían en qué clave
        corresponde a los mismos bytes: es la maleabilidad benigna
        de la figura~\ref{fig:param-narrowly}.
  \item Mensajería: \texttt{mceliece6960119} (el criptograma más
        pequeño entre los \texttt{6*}: $194$~B; \texttt{pc} si se
        quieren las defensas de fallo, $+32$~B): su clave pública de
        $\approx1$~MB se cachea una vez y se amortiza entre millones
        de mensajes. IoT: \texttt{mceliece460896f} (o
        \texttt{6960119f} si se exige largo plazo): la variante
        \texttt{f} porque la generación con reintentos es el punto
        débil de un dispositivo lento, y el tamaño de la clave
        pública allí importa menos que poder generarla de una vez.
\end{enumerate}

\subsection*{Capítulo 13 --- ¿Qué significa romper McEliece?
(capítulo~\ref{ch:seguridad})}

\begin{enumerate}[leftmargin=1.9em,itemsep=4pt]
  \item (a)~Con una clave equivalente ($g$ y el orden $\alpha$) se
        decodifica el código público: todo criptograma se abre, y
        con él todo mensaje. (b)~Un rompemensajes «de juego» triunfa
        sobre la \emph{distribución} del reto (por ejemplo con
        conocimiento del reparto uniforme de errores) y no tiene por
        qué resolver el problema de decodificación sobre síndromes
        arbitrarios: no hay reducción conocida de juego a
        decodificador general, y hay instancias estructuradas
        débiles. (c)~Un distinguidor solo decide «¿Goppa o
        aleatoria?»: no descifra; su papel es erosionar el supuesto
        de que la clave pública es pseudorandom (el modelo que usan
        las reducciones) y, a veces, ser la primera etapa de la
        recuperación estructural; la reducción
        distinguidor$\to$rotura es justamente la que no se conoce.
  \item Las $137$ palabras de peso $\le2$ y el síndrome
        $(\hx2,\hx b)$: el único candidato es $e_5+e_9$, porque
        $\tilde{\Hmat}e_5+\tilde{\Hmat}e_9
        =(\hx8,\hx e)+(\hx a,\hx5)=(\hx2,\hx b)$ (la unicidad estaba
        garantizada por $d=5$: dos candidatos diferirían en un
        codeword de peso $\le4<d$). Enumerar las palabras de peso
        $64$ de \texttt{348864} cuesta $\approx2^{456}$; Prange se
        cotiza en $2^{173{,}4}$: la diferencia ($\approx283$ bits)
        vive en el algoritmo: la estructura lineal convierte el
        «gran catálogo» en una probabilidad por iteración con
        iteraciones baratas (el acortamiento del
        capítulo~\ref{ch:isd}).
  \item La conjetura, con cuantificadores: \emph{para todo}
        algoritmo $A$ de tiempo subexponencial, \emph{para} la
        instancia $(H,\,C=He)$ generada por la distribución del
        esquema ($e$ uniforme de peso $t$), \emph{con éxito}
        $=\Pr[A(H,C)=e]$ despreciable (con los mejores ataques
        conocidos $\ge2^{140}$--$2^{150}$ en los modelos serios).
        No equivalen: (a)~«NP-duro» habla del peor caso,
        asintótico y de instancias arbitrarias: ni media, ni
        distribución, ni parámetros; (b)~«no existe algoritmo
        polinómico» ni fija la distribución ni cuantifica el éxito
        (y no implica la dureza concreta de las instancias del
        esquema); (c)~«el mejor ataque cuesta $2^{140{,}8}$» es una
        cota \emph{superior} empírica: evidencia, no teorema, y
        puede mejorar (2025--2026 lo mostró). En la práctica, el
        que sostiene las decisiones es (c): las categorías y los
        márgenes se calculan con el mejor ataque conocido; la
        conjetura es el marco de la promesa.
  \item $2^{140{,}8}$ son operaciones \emph{vectoriales} en el
        modelo de memoria libre, donde una operación de vector
        grande cuenta como una: no es directamente comparable con
        las $2^{143}$ operaciones de bits de buscar una clave
        AES-128; con el cargo realista de la memoria ($2^{86{,}6}$),
        el mismo ataque sube a $2^{158{,}6}$: factor
        $2^{158{,}6-140{,}8}=2^{17{,}8}\approx228\,210$. La
        categoría~$1$ se certifica con el modelo que carga la
        memoria ($158{,}6\ge143$): de ahí que no haya
        contradicción.
  \item Mitad en el exponente: $151{,}4/2=75{,}7$, muy por debajo
        del listón $2^{128}$: \texttt{348864} no pasa el test
        cuántico ingenuo. ISO estandarizó igualmente
        \texttt{460896} ($96{,}7$, también por debajo) porque la
        raíz cuadrada sobreestima al atacante: Grover está limitada
        por la latencia del bucle interno, y el coste cuántico real
        es mayor; con esa salvedad documentada, ISO prefirió
        incluir el conjunto de claves pequeñas antes que dejar el
        perfil sin cobertura.
  \item (a)~Distinguidor (familia estructural, capítulos
        \ref{ch:seguridad} y~\ref{ch:estructura}): rompe supuestos
        de indistinguibilidad; no descifra. (b)~Mejora de ISD
        (capítulo~\ref{ch:isd}): mueve la frontera de la
        decodificación genérica; sola no rompe el esquema: cambia
        la estimación. (c)~Canal lateral (capítulos
        \ref{ch:kem}--\ref{ch:parametros} y~\ref{ch:2026}): rompe
        un despliegue real (recupera la clave de una instalación
        concreta midiendo tiempos), no el esquema abstracto.
\end{enumerate}

\subsection*{Capítulo 14 --- Decodificación por conjuntos de
información (capítulo~\ref{ch:isd})}

\begin{enumerate}[leftmargin=1.9em,itemsep=4pt]
  \item (a)~$64\log_2(768/3488)=-139{,}73$;
        (b)~$k^3=2720^3$: $3\log_2 2720=34{,}23$; total
        $173{,}96$; (c)~la guía cotiza $173{,}4$: el medio bit son
        las constantes finas del estimador (trabajo de la
        eliminación, pivotaje, normalización del éxito) que
        nuestro producto grueso sobrecuenta. Diferencia de
        constantes, no de modelo.
  \item Primer salto: $22{,}0$ bits $=2^{22}\approx4{,}19$ millones
        $\approx4$ factores de $2^{20}$; segundo salto: $10{,}6$
        bits $\approx1552$, medio factor de $2^{20}$. El primero es
        el grande: el cruce de listas de Stern aporta
        $\approx4\cdot2^{20}$ sobre Prange; la búsqueda de vecinos
        de May--Ozerov añade «solo» $1{,}5\cdot2^{10}$.
  \item Porque ambos son $\bigl(1/(1-R)\bigr)^{t}$ por
        $2^{o(t)}$: con $R\approx0{,}78$, la base es
        $1/(1-R)\approx4{,}55$ ($2^{2{,}18t}$; para $t=64$,
        $2^{139{,}8}$). Prange arrastra un coste polinómico
        ($k^3=2^{34{,}2}$): de $173{,}96$; May--Ozerov mete el
        trabajo en las listas y los vecinos y se queda cerca del
        término desnudo ($140{,}8$). El «$5^t$» es taquigrafía del
        término dominante ($5\approx1/(1-R)$ redondeado): mismo
        exponente, $o(1)$ distinto, y a $t=64$ ese $o(1)$ vale
        $33$ bits.
  \item (a)~Memoria gratis vs realista para \texttt{6960119}:
        $279{,}2-245{,}7=33{,}5$ bits; (b)~$245{,}7<272$ (el
        ataque con memoria gratis sale «más barato» que agotar
        AES-256), pero $279{,}2>272$ (con el cargo de la memoria
        ya no): la categoría del equipo es la $5$ porque certifica
        con el presupuesto realista.
  \item (8192,96) se transporta a (4608,96) por acortamiento: un
        decodificador del código grande, restringido a las primeras
        $4608$ coordenadas, decodifica el pequeño; el éxito se
        conserva salvo el factor $\binom{4608}{96}/\binom{8192}{96}
        \approx2^{-0{,}8}$, que es la pérdida de apretamiento del
        transporte (ese $0{,}8$ es el exponente que separa ambas
        cifras en el modelo 1/2). La dirección inversa no se
        conoce: romper el pequeño no rompe automáticamente el
        grande.
  \item La mitad del exponente es una cota optimista porque
        supone que cada iteración de Grover cuesta lo mismo que su
        versión clásica y que el paralelismo no tiene penalización:
        en la práctica el bucle interno (accesos a memoria, la
        búsqueda en listas) domina y la ventaja cuántica real es
        menor que la cuadrática. Los $256$ bits de $\delta$ fijan
        el suelo complementario: la búsqueda cuántica de la semilla
        cuesta $2^{128}$ (Grover sobre $256$ bits, con una
        generación por comprobación), de modo que la seguridad
        cuántica de un conjunto es $\min($mitad-ISD$,2^{128})$: por
        eso el listón de ISO es exactamente $2^{128}$ y por eso los
        conjuntos por debajo ($75{,}7$, $96{,}7$) son el punto
        débil, aunque el modelo de la mitad sea optimista.
\end{enumerate}

\subsection*{Capítulo 15 --- Ataques estructurales: recuperar la
clave privada (capítulo~\ref{ch:estructura})}

\begin{enumerate}[leftmargin=1.9em,itemsep=4pt]
  \item (a)~$4096^{64}/64=2^{768-6}=2^{762}$;
        (b)~$1024^{50}/50\approx2^{500-5{,}6}=2^{494{,}4}$;
        (c)~$\log_2P(4096,769)=\sum_{j=0}^{768}\log_2(4096-j)
        \approx9116{,}8$.
  \item $\binom{4096}{3488}$ cuenta \emph{conjuntos} de $n$
        posiciones: $\approx2^{2476}$ (entropía de «qué elementos
        forman el soporte»); $P(4096,769)$ cuenta \emph{tuplas
        ordenadas} de $769$ posiciones concretas: el dato parcial
        que el modelo de 2022 necesitaba reconstruir. La diferencia
        ($\approx6640$ bits) es el precio del orden y de las
        posiciones concretas: una va con el tamaño del soporte
        ($n$) y la otra con los elementos libres ($n-k+1$).
  \item (a)~$9116{,}8-140{,}8=8976$ bits; el cociente de costes es
        $2^{8976}$ (¡$2^{8720}$ veces el margen entero de
        AES-256!). (b)~Comparada con $2^{256}$, la cifra de 2022 es
        inalcanzable por cualquier medio físico imaginable: el
        titular era «correcto» dentro de un modelo de datos
        parciales que estaba a punto de agotarse, y como amenaza
        práctica era casi cómico (nadie ejecutaría jamás tal
        cómputo).
  \item La cadena sería: distinguidor con ventaja $1-o(1)$
        $\to$ estadístico que expone estructura en las claves
        reales $\to$ (eslabón que falta) inversión eficiente en
        recuperación de clave $\to$ descifrado. Los márgenes de
        la tabla~\ref{tab:est-margenes} miden la distancia del régimen de
        cada método a los parámetros reales: Faugère ($t\le20$
        frente a $64$--$128$) tiene holgura enorme; Carrier
        (tasa $<0{,}3$ frente a $0{,}7$) moderada; el
        \emph{holdout}/anulación de orden superior ya alcanza el
        $t$ real: es el candidato natural a escalar, y el eslabón
        que falta es la barrera que hoy separa el estadístico de
        la clave.
  \item $n<q$ es doblemente bueno: contra ISD no se paga nada
        (el acortamiento transporta la dureza con $0{,}8$ bits de
        pérdida: observación~\ref{obs:isd-acortamiento}), y contra
        la división del soporte el atacante debe adivinar qué $n$
        de las $q$ columnas forman $L$ y en qué orden: paga
        $\binom{q}{n}$ (y tuplas ordenadas: $2^{9116{,}8}$); con
        $n=q$ ese frente se desvanece
        (observación~\ref{obs:est-nq}). El diseño cobra en un
        frente lo que no pierde en el otro.
  \item Holguras (menor = más peligroso): (1)~anulación de orden
        superior / \emph{holdout} (su régimen alcanza el $t$ real:
        la frontera está pegada a los parámetros, aunque sus
        costes aún no escalen); (2)~Carrier (necesita tasa baja,
        $<0{,}3$, y los conjuntos reales van a $0{,}7$--$0{,}8$:
        margen moderado en el otro lado de la frontera);
        (3)~Faugère ($t\le20$ frente a $\ge64$: holgura enorme).
        El más peligroso a escalar es el primero: es el único
        cuyo \emph{régimen} ya contiene a los conjuntos reales.
        (Ejercicio de razonamiento: la respuesta admite matices,
        pero no intercambiar los tres regímenes.)
\end{enumerate}

\subsection*{Capítulo 16 --- El terremoto: 2025--2026
(capítulo~\ref{ch:2026})}

\begin{enumerate}[leftmargin=1.9em,itemsep=4pt]
  \item (a)~$9116{,}8-94{,}3=9022{,}5$ bits de diferencia. No es el
        «mismo ataque mejorado»: la familia de 2022 (datos
        parciales y reconstrucción del soporte) estaba agotada
        contra su muro; el \emph{holdout} (GIJS/Weis) es una
        familia nueva, importada de la teoría de la información
        privada: no busca el soporte, evalúa el código con familias
        de funciones correlacionadas. (b)~$94{,}3<256$: por debajo
        del suelo trivial (enumerar la semilla de $256$ bits, con
        una generación por candidato). Una estimación así solo es
        compatible con (i)~una contabilidad que no cobra el trabajo
        real por candidato/paso (memoria gratis, «operaciones de
        bits», heurísticas no ejecutadas) o (ii)~una rotura
        histórica. De ahí la exigencia de verificación pública: a
        octubre de 2026, ninguna de las dos cosas está demostrada.
  \item Total con $100$--$1400$ ejecuciones: sin memoria,
        $2^{89}\cdot[100,1400]=2^{95{,}6}$--$2^{99{,}5}$; con cargo
        severo, $2^{117}\cdot[100,1400]=2^{123{,}6}$--$2^{127{,}5}$.
        La decodificación genérica en ese modelo para
        \texttt{348864} está en $2^{158{,}6}$: ambos rangos quedan
        por debajo: «nuestra mejor estimación; con el cargo de la
        guía queda por debajo de la decodificación genérica en ese
        modelo». Es exactamente la afirmación del equipo, y la
        razón del sobresalto: incluso el rango con cargo severo
        ($\le2^{127{,}5}$) roza el listón ISO de $2^{128}$.
  \item $94{,}3+8\cdot6=142{,}3\approx2^{142}$: la regla de Weis
        (unos $8$ bits por unidad de $m$) reproduce la instancia de
        Bernstein ($n=6176$, $t=112$, $m=18$), cifrada en
        $\approx2^{142}$, con error de $0{,}3$ bits. Clave pública
        allí: $mt\lceil k/8\rceil=2016\cdot520=\num{1048320}$~B
        $\approx1{,}0$~MB: cuatro veces la de \texttt{348864}
        ($0{,}25$~MB) y del orden de la mayor actual
        (\texttt{8192128}: $1{,}36$~MB). El eje $m$ compra
        seguridad exponencial a precio lineal en la clave.
  \item (a)~Demostrada (en su régimen asintótico: es un teorema
        con ventaja $1-o(1)$). (b)~Estimada: el método de
        recuperación es heurístico (cinco supuestos enumerados por
        Weis, probados solo en claves pequeñas y retos TII) y la
        cifra concreta $2^{102{,}4}$ depende de la contabilidad.
        (c)~Heurística y disputada: la afirmación de Vedenev
        «esencialmente una ejecución» no está verificada de forma
        independiente (hay réplicas y contra-réplicas). (d)~Hecho
        (negativo y verificable): nadie ha ejecutado el cómputo de
        núcleos contra un conjunto real; los retos TII tienen
        $n=214$.
  \item $2^{94{,}3}$ frente a $2^{143}$ (AES-128): en el mismo
        registro de «operaciones de bits», el ataque saldría
        $48{,}7$ bits más barato: alarmante. Pero ese registro no
        cobra la memoria ni la latencia; con el cargo severo
        ($2^{117}$--$2^{128}$), el ataque cuesta tanto o más que
        agotar AES-128: favorable al esquema. Falta: (i)~un modelo
        de costes consensuado (cómo cobrar memoria y accesos) y
        (ii)~una ejecución de extremo a extremo (nadie ha corrido
        el método contra un conjunto real): sin ambas, la
        comparación está abierta (BSI elige la precaución; el
        equipo, su contabilidad clásica).
  \item Ejemplos (límite de $80$ caracteres):
        \emph{verdadero} ---~«McEliece: ataques de 2026 estimados,
        ninguno ejecutado ni verificado»; \emph{falso por
        alarma}~--- «McEliece roto: un ataque de $2^{94}$ desarma
        el estándar ISO» (le falta: estimación, modelo de costes
        disputado, no ejecutado); \emph{falso por
        complacencia}~--- «McEliece intacto: nada ha cambiado»
        (le falta: las estimaciones bajaron $9000$ bits en cuatro
        años y el BSI desaconseja despliegues nuevos). (Ejercicio
        abierto: intercambiar con un colega y discutir la cláusula
        ausente.)
\end{enumerate}
